% ISR Framework
\section{Inflationary Sector Renormalization}
\label{sec:framework}

\subsection{Sector state and observable patch}

Let $x_i$ denote the state of the $i$th inflationary sector, described by
\begin{itemize}
  \item $\phi_i$ --- the final field value,
  \item $v_i$ --- the vacuum basin label,
  \item $\boldsymbol{\theta}(x_i) = (\Lambda_{\rm eff}, G_{\rm eff}, Q_{\rm fluct}, T_{\rm reh})_i$
        --- a vector of effective physics parameters assigned to the basin.
\end{itemize}
The observable patch $x_0$ is anchored at $\phi_0 = 2.0$ and its basin is
identified from the landscape potential, not from the first sampled trajectory.

\subsection{Locality kernels}

A kernel $K_\ell[d(x_i, x_0)]$ weights sectors by their distance to $x_0$.
We define three distance types:

\begin{itemize}
  \item \textbf{Field-space distance}: $d_\phi(x_i, x_0) = |\phi_i - \phi_0|$.
  \item \textbf{Initial-field distance}: $d_{\phi_{\rm init}}(x_i, x_0) = |\phi_{{\rm init},i} - \phi_0|$.
  \item \textbf{Basin-transition distance}: $d_{\rm basin}(x_i, x_0) = 0$ if
        $v_i = v_0$, and $1$ otherwise.
\end{itemize}

We use two evidence kernels:

\begin{itemize}
  \item Gaussian: $K_\ell(d) = \exp(-d^2 / 2\ell^2)$.
  \item Exponential: $K_\ell(d) = \exp(-|d| / \ell)$.
\end{itemize}

The scale parameter $\ell$ controls how quickly weight decays:
small $\ell$ isolates the local basin; large $\ell$ admits broader
unobserved-sector mixing.

\subsection{ISR measure}

The \ISR probability of observing effective physics $\boldsymbol{\theta}$ is

\begin{equation}
\label{eq:isr_measure}
P_{\rm ISR}(\boldsymbol{\theta} \mid x_0)
= \frac{\sum_i K_\ell[d(x_i,x_0)] \; \mathbb{1}[\boldsymbol{\theta}_i = \boldsymbol{\theta}]}
        {\sum_i K_\ell[d(x_i,x_0)]}.
\end{equation}

For comparison, the global census measure is

\begin{equation}
\label{eq:global_measure}
P_{\rm global}(\boldsymbol{\theta})
= \frac{1}{N} \sum_i \mathbb{1}[\boldsymbol{\theta}_i = \boldsymbol{\theta}].
\end{equation}

\subsection{Phenomenological status}
\label{sec:phenomenological_status}

\ISR is a conditional weighting scheme on a generated sector ensemble.
It does not regulate spacetime volume divergences and does not provide a
first-principles probability measure over vacua in an eternally inflating
multiverse.
Its purpose is to test whether a locality prior over unobserved
sectors---weighting field-space-nearby sectors more heavily---produces a
more cutoff-stable probability estimate within a finite trajectory sample
than naive equal weighting.

The toy model in this paper is a numerical testbed, not a realistic cosmology.
\ISR should be judged by stability, robustness, and explicit failure reporting,
not by whether it can be tuned to prefer a particular vacuum.
A measure that requires fine-tuned kernels and yields brittle results would be
uninteresting; a measure that transparently reports when and why it fails is
valuable even if imperfect.

\subsection{Kernel-scale selection rule}
\label{sec:kernel_scale_rule}

The kernel scale $\ell$ must be chosen large enough that the kernel
admits sectors from multiple basins.
If $\ell$ is much smaller than the inter-vacuum spacing in field space,
$N_{\rm eff} \ll N$ and the ISR distribution is effectively a delta function
on the $x_0$ basin, trivially suppressing KL drift.
This regime is a diagnostic artifact, not evidence of physical stability.
We do not require $\ell$ to exceed the full inter-vacuum spacing.
Instead, we define the non-degenerate regime empirically by requiring
$N_{\rm eff}/N \ge 0.3$, basin-exception rate $\le 0.3$, and performance
better than the shuffled-distance null.
In this landscape that selects approximately $0.5 \le \ell \le 1.5$,
with $\ell = 1.0$ used as the representative value.

\subsection{Cutoff stability diagnostic}

We assess measure stability by computing the KL divergence between
distributions at successive cutoffs $N_k$:

\begin{equation}
\label{eq:kl_diagnostic}
D_{\rm KL}(P_{N_k} \parallel P_{N_{k-1}})
= \sum_v P_{N_k}(v) \log\frac{P_{N_k}(v)}{P_{N_{k-1}}(v)}.
\end{equation}

We also compute the 1-Wasserstein distance over the cumulative
distribution function of vacuum probabilities.

A measure is considered stable if these divergences approach zero
as the cutoff increases.

\subsection{Kernel-weighted sector averaging}

Kernel-weighted averaging over unobserved sectors defines illustrative
aggregate coefficients for the observable patch:

\begin{equation}
\label{eq:effective_flow}
c_{\rm eff}(\ell) = \frac{\sum_i K_\ell[d(x_i,x_0)] \; c_i}
                         {\sum_i K_\ell[d(x_i,x_0)]},
\end{equation}

where $c_i$ is any component of $\boldsymbol{\theta}(x_i)$.
This mirrors a kernel-weighted coarse-graining: the weighting averages
over unobserved inflationary sectors rather than high-momentum modes.
